\documentclass[10pt,a4paper]{amsart}
\usepackage[margin=19mm]{geometry}
\usepackage{amsmath,amssymb,mathtools}
\usepackage[scr=rsfs,cal=euler]{mathalfa}
\usepackage{xfrac}
\usepackage[unicode,hidelinks,bookmarks=false]{hyperref}
\newcommand{\ie}{i.e.\ }
\newcommand{\cf}{cf.\ }
\newcommand{\resp}{respectively\ }
\newcommand{\Subsection}{\subsection}
\providecommand{\nobreakdash}{\nobreak\hbox{-}}
\setlength{\emergencystretch}{2em}
\multlinegap0pt
\usepackage{amssymb}
\usepackage{xcolor}
\usepackage{tikz}
\newtheorem{proposition}{Proposition}
\newcommand{\norm}[1]{\lVert{#1}\rVert}
\newcommand{\vtr}[1]{\mathbf{#1}}
\title{Property-Preserving Hashing for \texorpdfstring{$\ell_1$}{l1}-Distance Predicates: Applications to Countering Adversarial Input Attacks}
\author{Hassan Asghar, Chenhan Zhang, Dali Kaafar}
\date{Source: arXiv:2504.16355v2 (CC BY 4.0)}
\begin{document}
\maketitle

\noindent\emph{Source and licence.} Property-Preserving Hashing for \texorpdfstring{$\ell_1$}{l1}-Distance Predicates: Applications to Countering Adversarial Input Attacks. Version arXiv:2504.16355v2 (CC BY 4.0), distributed by arXiv under the Creative Commons Attribution 4.0 International licence, http://creativecommons.org/licenses/by/4.0/, as recorded in the arXiv licence metadata for this exact version and shown on the abstract page https://arxiv.org/abs/2504.16355v2. Full licence notice: \texttt{licenses/arxiv-2504.16355-CC-BY-4.0.txt}.\par\smallskip

\noindent\textit{Editorial scope.} Counted unit: the article's proposition prop:l1-l2-relation (source0.tex lines 214--222). The article's complete original proof of it is delivered with it (source0.tex lines 223--247, one \texttt{proof} environment of 25 lines). No mathematics is written, repaired or re-derived: the statement and the proof are the article's own lines, in the article's own notation, and no retained line is substituted or re-flowed.  No further source-local context is needed: every cross-reference of every retained line resolves inside the two delivered spans.  Nothing was omitted from any retained span, so the package carries no omission record.  No retained line contains a citation, so the wrapper carries no bibliography and no bibliography entry.  No retained line contains a figure environment, an image-inclusion command or drawing code, so no image is delivered and the package declares no asset dependency.  Editorial formatting only, in addition: an AMS \texttt{amsart} preamble that restates the article's own declarations and macros used by the retained lines, namely the environments \texttt{proposition} on separate counters, together with the standard packages those lines need.  The retained environments print in this document as Proposition 1 in order of appearance, whereas the article prints the same items with the numbering of its own full document.  The complete native source of the article as distributed in the arXiv e-print for this exact version is bundled unchanged as \texttt{sources/176940/source0.tex}.
\begin{proposition}
\label{prop:l1-l2-relation}
Let $\vtr{x}$ and $\vtr{y}$ be images. Then,
\begin{equation}
\label{eq:l1-l2-relation}
 \norm{\vtr{x} - \vtr{y}}_2 \leq \norm{\vtr{x} - \vtr{y}}_1 \leq \sqrt{n}\norm{\vtr{x} - \vtr{y}}_2.   
\end{equation}
Furthermore, these bounds are tight. 
\end{proposition}

\begin{proof}
Let $\vtr{z} = \vtr{x} - \vtr{y}$. Then, equivalently, we need to show that
\[
\norm{\vtr{z}}_2 \leq \norm{\vtr{z}}_1 \leq \sqrt{n}\norm{\vtr{z}}_2
\]
Note that $\vtr{z}$ may not be an image, i.e., $\vtr{z}$ may not be a member of $\mathbb{Z}_q^n$. Now we see that
\begin{align*}
  \norm{\vtr{z}}_1^2 &= \left( \sum_{i = 1}^n |z_i| \right)^2 \\
  &= \left( \sum_{i = 1}^n |z_i| \right) \left( \sum_{i = 1}^n |z_i| \right)\\
  &=  \sum_{i = 1}^n |z_i|^2  +  \sum_{i, j : i \neq j} |z_i||z_j| \\
  &\geq \sum_{i = 1}^n |z_i|^2 =  \norm{\vtr{z}}_2^2,
\end{align*}
from which it follows that $\norm{\vtr{z}}_2 \leq \norm{\vtr{z}}_1$. For the second inequality, let $\vtr{1}$ be the vector of all 1's, and let $\vtr{b}$ be such that $b_i = |z_i|$ for all $i$. Then,
\begin{equation*}
  \norm{\vtr{z}}_1 = \sum_{i=1}^n |z_i| = \langle \vtr{1}, \vtr{b} \rangle  \leq \norm{\vtr{1}}_2 \norm{\vtr{b}}_2 = \sqrt{n} \norm{\vtr{z}}_2,
\end{equation*}
where we have used the Cauchy-Schwarz inequality. To show that the bounds are tight, let us first consider the inequality on the left. Assume that for some positive constant $c$ we have $c \norm{\vtr{x} - \vtr{y}}_2 \leq \norm{\vtr{x} - \vtr{y}}_1$. Consider the two images $\vtr{x} = (1, 0, \ldots, 0)$ and $\vtr{y} = \vtr{0}$, where $\vtr{0}$ is the zero vector. Then,
\[
c \norm{\vtr{x} - \vtr{y}}_2 \leq \norm{\vtr{x} - \vtr{y}}_1 \Rightarrow c \sqrt{1} \leq 1 \Rightarrow c \leq 1.
\]
Thus, $c = 1$ is tight. For the RHS, assume that for some positive constant $c$, we have $ \norm{\vtr{x} - \vtr{y}}_1 \leq c \norm{\vtr{x} - \vtr{y}}_2$.  Consider the two images $\vtr{x} = (q-1, q-1, \ldots, q-1)$ and $\vtr{y} = \vtr{0}$. Then,
\[
\norm{\vtr{x} - \vtr{y}}_1 \leq c \norm{\vtr{x} - \vtr{y}}_2 \Rightarrow n(q-1) \leq c \sqrt{n}(q-1) \Rightarrow c \geq \sqrt{n}. 
\]
\end{proof}

\end{document}
